\par\clearpage
\noindent{\small\textbf{ZUR HANNOVERSCHEN TRIANGULATION.}}\hfill359\par\medskip
\begin{center}
\textsc{Gauss} an \textsc{Bessel}.\quad Göttingen, 5.~November 1823.
\end{center}
\begingroup
\makeatletter
\renewcommand{\@makefnmark}{\hbox{\normalfont[\textsuperscript{*})]}}
\long\def\@makefntext#1{\parindent1.5em\noindent[\textsuperscript{*})\ #1}
\makeatother

\par\noindent\hspace*{\parindent}.\ .\ .\ .\ .\ .\quad Ich habe einen Theil des Jahrs damit zugebracht, den noch übrigen Theil meiner trigonometrischen Messungen im Norden: von Timpenberg, Nindorf, Lüneburg bis Hamburg zu absolviren; dann auch vorläufig die weiter westlich liegende Gegend, nach Bremen zu, zu recognosciren, da unser Gouvernement eine weitere Fortsetzung der Messungen nach Westen zu, bis zur holländischen Grenze und zum Anschluss an die \textsc{Krayenhoff}schen Dreiecke, wünscht; endlich zuletzt habe ich noch einmal den Brocken und Hohehagen besucht, da theils die Winkel des Dreiecks Hohehagen-Hils-Brocken im Jahr 1821 unter sehr ungünstigen Umständen gemessen, theils jetzt noch die damals missglückte Verbindung des Brockens mit dem Inselsberg zu effectuiren war, sowie auch jetzt noch der hessische Dreieckspunkt Meisner angeknüpft werden sollte.
Ich habe diese Zwecke \textls[120]{meistens} zu meiner Zufriedenheit erreicht; nur den Hils hätte ich gern auch noch einmal besucht, um den Winkel dort genauer zu messen; die vorgerückte Jahreszeit hat mich aber daran gehindert.
Diesen Umstand abgerechnet, kann ich jetzt die Triangulirung zur Gradmessung, so weit sie zu meinem Ressort gehört, als geendigt ansehen.
Astronomische Beobachtungen sind noch keine weiter gemacht, als die zur Orientirung meiner ersten Dreiecksseite gehören.
Ob \textls[120]{ich} den oben erwähnten Plan der Fortsetzung der Messungen nach Westen noch ausführe, ist übrigens noch sehr ungewiss.
Es ist manches dafür, manches, fast noch mehr, dagegen, auch abgesehen davon, dass vielleicht noch die Möglichkeit einer Änderung meiner äussern Lage eintreten könnte\footnote{Es handelte sich um die Berufung nach Berlin.]}.\par

\par\noindent\hspace*{\parindent}.\ .\ .\ .\ .\ .\quad Ich habe das System meiner Hauptdreiecke in diesen Tagen sorgfältig ausgeglichen, so dass nicht nur die Summe der Winkel jedes einzelnen Dreiecks, sondern auch die Verhältnisse der Seiten in den gekreuzten Vierecken und Fünfecken genau harmoniren, und zwar ohne alle Willkür, ohne Auswählen, ohne Ausschliessen, alles nach der Strenge der Probabilitätsrechnung.
Es sind zusammen 26 Dreiecke, worin alle Winkel von mir selbst beobachtet sind.
\par\clearpage
\noindent360\hfill{\small\textbf{BRIEFWECHSEL MIT BESSEL.}}\par\medskip
\noindent Die grösste Summe der Fehler ist \(2^{\prime\prime}{,}2\) in einem Dreiecke, wo bei einer Seite das Pointiren sehr schwierig war; die nächst grösste ist \(1^{\prime\prime}{,}8\).
Keine der 76 vorkommenden Richtungen ist bei der Ausgleichung um eine ganze Secunde geändert; die grösste Änderung beträgt \(0^{\prime\prime}{,}813\) bei der oben erwähnten Seite von Nindorf nach Hamburg.
Was ich nach meiner neuen Probabilitätstheorie den mittlern Fehler nenne, bei den Richtungen, ist \(0^{\prime\prime}{,}48\).
\par\endgroup

